%% toc: P3 40 40.8
\section{Proofs for nonconvergence and the limits of parity definability}
\label{sec:annex-p03-w03-p03-071-nonconvergence}

Throughout, the conventions of \cref{def:p03-w03-p03-071-assumptions} are in
force.  In \(G_<(n,p)\), for \(u<e\) the trace \(N_u(e)\) is determined by
the \(u\) pairs \(\{w,e\}\) with \(w\leq u\); distinct \(e>u\) use disjoint
sets of pairs.  Hence the traces \(N_u(e)\), \(u<e\leq n\), are mutually
independent, and for \(p=\frac12\) each is uniformly distributed over the
\(2^u\) subsets of \(\{1,\ldots,u\}\).  This fact is used repeatedly below
and is called the \emph{trace fact}.

\begin{annexitem}{annex:p03-w03-p03-071-no-zero-one}{thm:p03-w03-p03-071-no-zero-one}{No zero--one law with a built-in order}
\begin{proof}
Let \(n\geq2\).  The matrix of \(B\) requires \(x<y\) and that every
\(z<y\) equals \(x\); so exactly one element lies below \(y\), which forces
\(y=2\) and \(x=1\).  Conversely \(x=1,y=2\) satisfy the first three
conjuncts.  Hence \(G_<(n,p)\entails B\) exactly when \(1\sim2\), an event of
probability \(p\).  For \(p=\frac12\) the constant sequence \(\frac12\) tends
neither to \(0\) nor to \(1\).
\end{proof}
\end{annexitem}

\begin{annexitem}{annex:p03-w03-p03-071-unique}{thm:p03-w03-p03-071-unique}{The arithmetization is unique}
\begin{proof}
Let \(+,\ast\) satisfy items 1--4 of
\cref{def:p03-w03-p03-071-arithmetizable}.  We show by induction on
\(y\in S\) that for all \(x,z\in S\), \(+(x,y,z)\) holds if and only if
\(x+y=z\).  For \(y=1\) this is item~1.  Let \(y>1\).  By item~2 and the
induction hypothesis, \(+(x,y,z)\) holds if and only if \(z\neq1\) and
\(x+(y-1)=z-1\), that is, if and only if \(z\neq1\) and \(x+y=z\).  Since
\(x+y\geq2\), the condition \(z\neq1\) follows from \(x+y=z\), so
\(+(x,y,z)\) holds exactly when \(x+y=z\).  Thus \(+=\mathrm{Add}_u\).

Next, by induction on \(y\), \(\ast(x,y,z)\) holds if and only if \(xy=z\).
For \(y=1\) this is item~3.  Let \(y>1\).  By item~4, the induction
hypothesis and \(+=\mathrm{Add}_u\), \(\ast(x,y,z)\) holds if and only if
some \(v\in S\) has \(v=x(y-1)\) and \(x+v=z\).  If \(x(y-1)\leq u\) this says
\(z=x+x(y-1)=xy\).  If \(x(y-1)>u\) there is no such \(v\), and also no
\(z\in S\) with \(xy=z\), because \(xy>x(y-1)>u\).  In both cases
\(\ast(x,y,z)\) holds exactly when \(xy=z\), so \(\ast=\mathrm{Mul}_u\).

Conversely, \(\mathrm{Add}_u\) and \(\mathrm{Mul}_u\) satisfy items 1--4: item~1
and item~3 are the facts \(x+1=z\iff z=x+1\) and \(x\cdot1=x\); item~2
holds because for \(y>1\), \(x+y=z\) is equivalent to \(z\neq1\wedge
x+(y-1)=z-1\); item~4 holds because for \(y>1\), \(xy=z\) with \(z\in S\) is
equivalent to the existence of \(v=x(y-1)\in S\) with \(x+v=z\), as computed
above.
\end{proof}
\end{annexitem}

\begin{annexitem}{annex:p03-w03-p03-071-literal-fails}{thm:p03-w03-p03-071-literal-fails}{The literal encoding never represents addition}
\begin{proof}
Let \(u\geq2\), so \(1,2\in S\).  Suppose \(R_{c,d}(1,1,2)\).  Then some
\(e\) with \(c\leq e<d\) is adjacent to \(1\) and \(2\) and to no other
element of \(S\).  The triple \((1,2,2)\) has \(1,2\leq2\), and the same
\(e\) is adjacent to \(1,2,2\) and to no other element of \(S\); so
\(R_{c,d}(1,2,2)\).  Since \((1,1,2)\in\mathrm{Add}_u\) and
\((1,2,2)\notin\mathrm{Add}_u\), either \(R_{c,d}(1,1,2)\) fails or
\(R_{c,d}\) contains a non-sum; in both cases \(R_{c,d}\neq\mathrm{Add}_u\).

For multiplication, \((2,1,2)\in\mathrm{Mul}_u\).  If \(R_{c,d}(2,1,2)\), a
vertex \(e\in[c,d)\) is adjacent to exactly the elements \(1,2\) of \(S\),
and then \(R_{c,d}(1,1,2)\), since \(1,1\leq2\); but \(1\cdot1\neq2\).  So
\(R_{c,d}\neq\mathrm{Mul}_u\).

If \(S\) were first order arithmetizable in Spencer's sense, the relations
\(R_{c,d}\) and \(R_{c',d'}\) would arithmetize \(S\), and by
\cref{thm:p03-w03-p03-071-unique} \(R_{c,d}=\mathrm{Add}_u\), which was just
excluded.  Spencer's representation lemma asserts that, almost surely,
\(R_{c,d}\) is \(+\) on \(\{1,\ldots,u\}\) for some \(c\leq d\); for
\(u\geq2\) this has probability \(0\).
\end{proof}
\end{annexitem}

\begin{annexitem}{annex:p03-w03-p03-071-witness}{thm:p03-w03-p03-071-witness}{What a witness interval is}
\begin{proof}
For \(x,y,z,e\) with \(x,y,z\leq u\), \(\operatorname{Tr}_u(e;x,y,z)\) holds
exactly when \(x<z\), \(y<z\) and \(N_u(e)=\{x,y,z\}\): the three adjacencies
give \(\{x,y,z\}\subseteq N_u(e)\), and the universal clause gives the
reverse inclusion.  Hence
\[
  R^<_{c,d}=\{(x,y,z)\in S^3:x<z,\ y<z,\ \{x,y,z\}=N_u(e)\text{ for some }
  e\in[c,d)\}. \tag{\(\dagger\)}
\]
A set has the form \(\{x,y,z\}\) with \(x,y<z\) exactly when it has two
elements (take \(x=y\) the smaller, \(z\) the larger) or three (take \(z\)
the largest).

\emph{Item 1, addition.}  Every set of \(\mathcal F_+(u)\) has two or three
elements: \(\{a,2a\}\) when \(a=b\), and \(\{a,b,a+b\}\) with
\(a<b<a+b\) otherwise.  As \(c>u\) and the traces of \(c,\ldots,d-1\) are
these sets, \(\operatorname{Clean}_u(c,d)\) holds; for \(m_+(u)=0\) the
interval is empty and the clause is vacuous.  By \((\dagger)\),
\(R^<_{c,d}\) consists of the triples \((x,y,z)\) with \(x,y<z\) and
\(\{x,y,z\}\in\mathcal F_+(u)\).  If \(\{x,y,z\}=\{a,b,a+b\}\) with
\(a<b\), then \(z\), being larger than \(x\) and \(y\), is the largest
element \(a+b\), and \(\{x,y\}=\{a,b\}\); so \(x+y=z\).  If
\(\{x,y,z\}=\{a,2a\}\), then \(z=2a\) and \(x=y=a\); so \(x+y=z\).
Conversely, for \((x,y,x+y)\in\mathrm{Add}_u\) we have \(x,y<x+y\) and
\(\{x,y,x+y\}=\{a,b,a+b\}\) with \(a=\min(x,y)\), \(b=\max(x,y)\).  Hence
\(R^<_{c,d}=\mathrm{Add}_u\).

\emph{Item 1, multiplication.}  Every set of \(\mathcal F_\times(u)\) has two
or three elements: \(\{a,a^2\}\) with \(a<a^2\) when \(a=b\geq2\), and
\(\{a,b,ab\}\) with \(a<b<ab\) when \(2\leq a<b\).  So
\(\operatorname{Clean}_u(c',d')\) holds as before.  By \((\dagger)\) and the
same reading of the largest element, \(R^<_{c',d'}(x,y,z)\) gives
\(\{x,y\}=\{a,b\}\), \(z=ab\) or \(x=y=a\), \(z=a^2\); in both cases
\(xy=z\).  The unit disjuncts of \(T\) give \(x\cdot1=x\) and \(1\cdot y=y\).
So \(T\subseteq\mathrm{Mul}_u\).  Conversely let \(xy=z\) in \(S\).  If
\(y=1\) or \(x=1\) a unit disjunct holds.  Otherwise \(x,y\geq2\), so
\(x,y<xy=z\) and \(\{x,y,z\}\in\mathcal F_\times(u)\), whence
\(R^<_{c',d'}(x,y,z)\).  So \(T=\mathrm{Mul}_u\).

\emph{Item 2.}  \(\operatorname{Good}\) is the conjunction of the two
cleanliness clauses with items 1--4 of
\cref{def:p03-w03-p03-071-arithmetizable} for the relations \(P,T\subseteq
S^3\) defined by the displayed formulas.  By
\cref{thm:p03-w03-p03-071-unique}, items 1--4 hold for \(P,T\) exactly when
\(P=\mathrm{Add}_u\) and \(T=\mathrm{Mul}_u\).

\emph{Item 3.}  Let \(F\in\mathcal F_+(u)\), \(F=\{a,b,a+b\}\) with
\(a\leq b\).  Then \((a,b,a+b)\in\mathrm{Add}_u=R^<_{c,d}\), so by
\((\dagger)\) some \(e\in[c,d)\) has \(N_u(e)=F\).  Distinct sets of
\(\mathcal F_+(u)\) are traces of distinct vertices, so \([c,d)\) has at
least \(m_+(u)\) elements, and \(c,\ldots,c+m_+(u)-1\) belong to it.  These
lie above \(u\) because \(u<c\), and in \(\Omega\) because \([c,d)\subseteq
\Omega\).  By \(\operatorname{Clean}_u(c,d)\) each has a trace of the form
\(\{x,y,z\}\) with \(x,y<z\), which has two or three elements.
\end{proof}
\end{annexitem}

\begin{annexitem}{annex:p03-w03-p03-071-representation}{thm:p03-w03-p03-071-representation}{Representation lemma}
\begin{proof}
Let \(U=\lfloor0.9L\rfloor\) and \(M=\lfloor n/2\rfloor\); take \(n\) so large
that \(n/4\geq1+L+L^2\).  Fix \(u\) with \(1\leq u\leq U\).

\emph{Addition.}  Write \(m=m_+(u)\) and list
\(\mathcal F_+(u)=\{F_0,\ldots,F_{m-1}\}\).  The map
\((a,b)\mapsto\{a,b,a+b\}\) on pairs \(1\leq a\leq b\) with \(a+b\leq u\) is
injective, since the set determines \(a+b\) as its largest element and then
\(\{a,b\}\) as the rest (or \(a=b\) when only two elements remain); so
\(m\leq u^2\).  If \(m=0\), put \(c=d=u+1\); by
\cref{thm:p03-w03-p03-071-witness}, item~1, \(\operatorname{Clean}_u(c,d)\)
and \(R^<_{c,d}=\mathrm{Add}_u\) hold with certainty.  If \(m\geq1\), let
\(K=\lfloor(M-u)/m\rfloor\) and \(c_j=u+1+jm\) for \(0\leq j<K\).  The blocks
\([c_j,c_j+m)\) are disjoint and contained in \((u,M]\).  Let \(W_j\) be the
event that \(N_u(c_j+i)=F_i\) for all \(i<m\).  By the trace fact,
\(\Pr[W_j]=2^{-um}\) and the events \(W_j\) are independent.  So
\[
  \Pr\Bigl[\bigcap_{j<K}\neg W_j\Bigr]=(1-2^{-um})^K\leq\exp(-K2^{-um}).
\]
Now \(um\leq u^3\leq(0.9)^3\log n<0.73\log n\), so \(2^{-um}\geq n^{-0.73}\);
and \(K\geq(n/2-1-u)/u^2-1\geq n/(4L^2)\) by the choice of \(n\).  Hence
\(K2^{-um}\geq n^{0.27}/(4(\log n)^{2/3})\), which is at least \(n^{1/4}\)
for all large \(n\).  If some \(W_j\) occurs, then \(c=c_j\), \(d=c_j+m\)
satisfy \(\operatorname{Clean}_u(c,d)\) and \(R^<_{c,d}=\mathrm{Add}_u\) by
\cref{thm:p03-w03-p03-071-witness}, item~1.

\emph{Multiplication.}  The same argument with \(\mathcal F_\times(u)\),
\(m_\times(u)\leq u^2\), and blocks \(c'_j=M+1+jm_\times(u)\) inside
\((M,n]\), which also lies above \(u\), gives \(c',d'\) with
\(\operatorname{Clean}_u(c',d')\) and \(T=\mathrm{Mul}_u\), except on an event
of probability at most \(\exp(-n^{1/4})\) for large \(n\); for
\(m_\times(u)=0\) take \(c'=d'=u+1\).

By \cref{thm:p03-w03-p03-071-witness}, item~2, the chosen parameters satisfy
\(\operatorname{Good}(u,c,d,c',d')\), so \(\operatorname{Ar}(u)\).  The
probability that this fails for some \(u\leq U\) is at most
\[
 2U\exp(-n^{1/4})\leq2L\exp(-n^{1/4}),
\]
which tends to \(0\).
\end{proof}
\end{annexitem}

\begin{annexitem}{annex:p03-w03-p03-071-upper}{thm:p03-w03-p03-071-upper}{Upper limit of the encoding}
\begin{proof}
Let \(u\geq30\) and write \(m=m_+(u)\).  By the injectivity noted in the
proof of \cref{thm:p03-w03-p03-071-representation}, \(m\) is the number of
pairs \(1\leq a\leq b\) with \(a+b\leq u\).  For each \(z\) with
\(2\leq z\leq u\) there are \(\lfloor z/2\rfloor\geq(z-1)/2\) such pairs with
\(a+b=z\), so \(m\geq\sum_{z=2}^{u}(z-1)/2=u(u-1)/4\geq u^2/5\), the last step
because \(u\geq5\).

Suppose \(\operatorname{Ar}(u)\), witnessed by \(c,d,c',d'\).  By
\cref{thm:p03-w03-p03-071-witness}, items 2 and 3, the event \(E_{u,c}\)
occurs: \(u<c\), \(c+m-1\leq n\), and each of \(c,\ldots,c+m-1\) has a trace
on \(\{1,\ldots,u\}\) of two or three elements.  Let \(Q_u\) be the family
of such sets; \(|Q_u|=\binom u2+\binom u3\leq u^3\).  By the trace fact,
\(\Pr[E_{u,c}]=(|Q_u|2^{-u})^m\).  We have \(u^3\leq2^{u/2}\) for
\(u\geq30\): it holds at \(u=30\) (\(27000\leq32768\)), and from \(u\) to
\(u+1\) the left side grows by the factor \((1+1/u)^3\leq(31/30)^3<1.11\)
and the right by \(\sqrt2>1.41\).  So
\[
  \Pr[\operatorname{Ar}(u)]\leq\sum_{c}\Pr[E_{u,c}]
  \leq n\,(2^{-u/2})^{u^2/5}=n\,2^{-u^3/10}.
\]

For the second claim let \(u_0=\lceil3L\rceil\), and let \(n\) be so large
that \(u_0\geq30\).  Since \((u_0+j)^3\geq u_0^3+j\),
\[
  \Pr[\exists u\geq u_0:\operatorname{Ar}(u)]
  \leq n\sum_{j\geq0}2^{-(u_0^3+j)/10}
  =\frac{n\,2^{-u_0^3/10}}{1-2^{-1/10}}\leq15\,n\,2^{-2.7\log n}
  =15\,n^{-1.7},
\]
using \(u_0^3\geq27\log n\) and \(1-2^{-1/10}>1/15\).  This tends to \(0\).
\end{proof}
\end{annexitem}

\begin{annexitem}{annex:p03-w03-p03-071-log-mod}{thm:p03-w03-p03-071-log-mod}{The residue of \(\lfloor\log u\rfloor\) is definable}
\begin{proof}
All elements below lie in \(S=\{1,\ldots,u\}\), and \(u\geq2^{40}\).

(i) \(\operatorname{One}(x)\) holds exactly for \(x=1\).
\(\operatorname{Ev}(a)\) holds exactly for even \(a\): if \(a=2b\) then
\(b\leq a\), so \(b\in S\) and \(\mathrm{Add}(b,b,a)\); conversely
\(\mathrm{Add}(b,b,a)\) means \(a=2b\).

(ii) \(\operatorname{Pw}(w)\) holds exactly when \(w\) is a power of \(2\).
If \(w=2^j\) and \(\mathrm{Mul}(a,b,w)\), then \(a\) divides \(2^j\), so
\(a=2^i\), which is \(1\) or even.  If \(w\) is not a power of \(2\), it has
an odd prime divisor \(a\); then \(b=w/a\in S\), \(\mathrm{Mul}(a,b,w)\), and
\(a\) is neither \(1\) nor even.

(iii) So \(\operatorname{Top}(w)\) holds exactly for \(w=2^k\), the largest
power of \(2\) not above \(u\), where \(k=\lfloor\log u\rfloor\).

(iv) By induction on \(i\), \(\delta_i(x)\) holds exactly for \(x=2^i\) when
\(2^i\leq u\), and for no \(x\) otherwise: \(\delta_{i+1}(x)\) says that
\(x=y+y\) for the element \(y=2^i\).  Since \(2^{40}\leq u\), \(\kappa(x)\)
holds exactly for \(x=q:=2^{40}-1\).

(v) Let \(1\leq r\leq20\).  For \(w\in S\), \(\operatorname{Cong}_r(w)\)
holds exactly when \(w=tq+2^r\) for some integer \(t\geq0\): the case
\(t=0\) is \(\delta_r(w)\), and for \(t\geq1\) the numbers \(t\leq tq=v\leq
w\), \(q\) and \(h=2^r\) lie in \(S\), so the second disjunct expresses
\(w=tq+2^r\).

(vi) Let \(w=2^k\) and \(r'=k\bmod40\).  Because \(2^{40}\equiv1\pmod q\),
\(2^k\equiv2^{r'}\pmod q\).  If \(r'=r\), then \(k\geq r\), \(40\) divides
\(k-r\), and \(q=2^{40}-1\) divides \(2^{k-r}-1\); so
\(2^k-2^r=2^r(2^{k-r}-1)=tq\) with \(t\geq0\), and
\(\operatorname{Cong}_r(2^k)\).  If \(\operatorname{Cong}_r(2^k)\), then
\(2^r\equiv2^k\equiv2^{r'}\pmod q\); both \(2^r\) and \(2^{r'}\) lie in
\([1,2^{39}]\subseteq[0,q)\), so they are equal and \(r'=r\).

Together, \(\Phi\) holds exactly when \(k\bmod40=r\) for some
\(r\in\{1,\ldots,20\}\).
\end{proof}
\end{annexitem}

\begin{annexitem}{annex:p03-w03-p03-071-kary}{thm:p03-w03-p03-071-kary}{Every \(k\)-ary relation on a short initial segment is represented}
\begin{proof}
Fix \(k\geq1\), write \(U=U_k(n)\), and let \(n\) be so large that
\(\log n\geq1\) and \(n\geq3\log n+2\).  If \(U=0\) there is nothing to
prove, so let \(U\geq1\).  Then \(U^k\leq U^{k+1}\leq\log n/(2(k+1))\leq\log
n\).

\emph{The pattern for one relation.}  Fix \(u\) with \(1\leq u\leq U\) and
\(R\subseteq S^k\), \(S=\{1,\ldots,u\}\), and let \(r=|R|\leq u^k\).  If
\(r=0\), take \(c=d=u+1\), which lies in \(\Omega\) because \(u\leq\log
n<n\).  A witness \(e_0,\ldots,e_k\) of
\(\operatorname{Rep}^{(k)}_{u,c,d}\) would satisfy \(c\leq e_0\leq e_k<d=c\),
which is impossible; so \((c,d)\) represents \(\varnothing=R\).  Let
\(r\geq1\), list \(R=\{\bar a^{(0)},\ldots,\bar a^{(r-1)}\}\), and put
\(b=(k+1)r\).  The \emph{pattern} of \(R\) is the sequence
\(P_0,\ldots,P_{b-1}\) of subsets of \(S\) with \(P_{i(k+1)}=\varnothing\) and
\(P_{i(k+1)+m}=\{a^{(i)}_m\}\) for \(i<r\) and \(1\leq m\leq k\).

\emph{Correctness of a pattern.}  Suppose \(u<c\), \(d=c+b\leq n\) and
\(N_u(c+p)=P_p\) for every \(p<b\).  We show that \((c,d)\) represents \(R\).
For \(i<r\) the vertices \(e_m=c+i(k+1)+m\), \(0\leq m\leq k\), are
consecutive, satisfy \(c\leq e_0\) and \(e_k\leq c+b-1<d\), and have traces
\(N_u(e_0)=\varnothing\) and \(N_u(e_m)=\{a^{(i)}_m\}\); so they witness
\(\operatorname{Rep}^{(k)}_{u,c,d}(\bar a^{(i)})\).  Conversely let
\(\bar x\in S^k\) and let \(e_0,\ldots,e_k\) witness
\(\operatorname{Rep}^{(k)}_{u,c,d}(\bar x)\).  Then
\(c\leq e_0<\cdots<e_k<d\), so all lie in \(\{c,\ldots,c+b-1\}\).  Let
\(p=e_0-c\).  As \(N_u(e_0)=\varnothing\) and the only empty sets of the
pattern are those at positions divisible by \(k+1\) (the others are
singletons), \(p=i(k+1)\) with \(i<r\).  Then \(e_m=c+i(k+1)+m\) has trace
\(\{a^{(i)}_m\}\), and the formula requires the trace of \(e_m\) to be
\(\{x_m\}\); so \(x_m=a^{(i)}_m\) for every \(m\), and \(\bar x=\bar
a^{(i)}\in R\).

\emph{Probability.}  Let \(K=\lfloor(n-u-1)/b\rfloor\) and \(c_j=u+1+jb\) for
\(0\leq j<K\); then \(c_j+b\leq u+1+Kb\leq n\).  Let \(W_j\) be the event
\(N_u(c_j+p)=P_p\) for all \(p<b\).  By the trace fact the events \(W_j\)
are independent and each has probability \(2^{-ub}\).  Now
\(ub\leq(k+1)u^{k+1}\leq(k+1)U^{k+1}\leq\frac12\log n\), so
\(2^{-ub}\geq n^{-1/2}\); also \(b\leq(k+1)U^k\leq\frac12\log n\) and
\(u+1\leq\log n+1\), so
\(K\geq(n-\log n-1)/(\tfrac12\log n)-1\geq n/\log n\) by the choice of
\(n\): with \(\lambda=\log n\), the middle expression is
\((2n-3\lambda-2)/\lambda\), which is at least \(n/\lambda\) because
\(n\geq3\lambda+2\).  Hence
\[
  \Pr\Bigl[\bigcap_{j<K}\neg W_j\Bigr]=(1-2^{-ub})^K\leq\exp(-K2^{-ub})
  \leq\exp\bigl(-n^{1/2}/\log n\bigr).
\]
If some \(W_j\) occurs, \((c_j,c_j+b)\) represents \(R\), by the previous
paragraph.

\emph{All relations at once.}  The number of pairs \((u,R)\) with
\(1\leq u\leq U\) and \(R\subseteq\{1,\ldots,u\}^k\) is at most
\(U\,2^{U^k}\leq\log n\cdot2^{\log n}=n\log n\).  So the probability that
some such \(R\) is not represented is at most
\(n\log n\cdot\exp(-n^{1/2}/\log n)\), which tends to \(0\).
\end{proof}
\end{annexitem}

\begin{annexitem}{annex:p03-w03-p03-071-prime}{thm:p03-w03-p03-071-prime}{Primes and Fermat primes are definable}
\begin{proof}
All elements below lie in \(S=\{1,\ldots,u\}\).  By the proof of
\cref{thm:p03-w03-p03-071-log-mod}, items (i) and (ii), whose arguments do
not use \(u\geq2^{40}\), \(\operatorname{One}(x)\) holds exactly for \(x=1\)
and \(\operatorname{Pw}(w)\) exactly when \(w\) is a power of \(2\).
Clearly \(\operatorname{Div}(a,b)\) holds exactly when \(a\) divides \(b\),
since a quotient \(b/a\) is at most \(b\) and so lies in \(S\);
\(\operatorname{Two}(w)\) exactly when \(w=2\); and \(\operatorname{Sq}(w)\)
exactly when \(w\) is a perfect square, since a square root of \(w\) is at
most \(w\).  A power \(2^i\) is a perfect square exactly when \(i\) is even:
if \(y^2=2^i\), every prime factor of \(y\) divides \(2^i\), so \(y=2^a\) and
\(i=2a\).

\emph{Primes.}  If \(x\) is prime, then \(x\neq1\), and \(ab=x\) forces
\(a=1\) or \(b=1\).  If \(x\) is not prime, either \(x=1\), or \(x=ab\) with
\(1<a,b<x\), and then \(a,b\in S\) and \(\mathrm{Mul}(a,b,x)\) with neither
factor equal to \(1\).

\emph{A divisibility fact.}  For integers \(i,j\geq1\), \(2^i-1\) divides
\(2^j-1\) if and only if \(i\) divides \(j\).  If \(j=im\), then
\(2^j-1=(2^i)^m-1=(2^i-1)\sum_{l<m}2^{il}\).  Conversely write \(j=im+\varrho\)
with \(0\leq\varrho<i\); then
\(2^j-1=2^\varrho(2^{im}-1)+(2^\varrho-1)\).  If \(2^i-1\) divides
\(2^j-1\), it divides \(2^{im}-1\) by the first part, hence divides
\(2^\varrho-1\), which lies in \([0,2^i-1)\); so \(2^\varrho-1=0\) and
\(\varrho=0\).

\emph{Fermat primes, forward.}  Suppose \(\operatorname{Fer}(x)\) with
witnesses \(v,v'\).  Then \(x\) is prime, \(x=v+1\), \(v=v'+1\geq2\), and
\(v=2^j\) with \(j\geq1\).  Let \(i\) be a divisor of \(j\).  Then
\(w=2^i\leq v\) lies in \(S\), \(\operatorname{Pw}(w)\) holds,
\(w'=2^i-1\geq1\) satisfies \(w=w'+1\), and \(w'\) divides \(v'=2^j-1\) by
the divisibility fact; so \(\operatorname{Two}(w)\) or
\(\operatorname{Sq}(w)\), that is, \(i=1\) or \(i\) is even.  Write
\(j=2^mo\) with \(o\) odd.  Then \(o\) divides \(j\) and is odd, so \(o=1\).
Hence \(x=2^{2^m}+1\) is a Fermat prime.

\emph{Fermat primes, backward.}  Let \(x=2^{j}+1\) be prime with
\(j=2^m\).  Put \(v=2^j\geq2\) and \(v'=v-1\geq1\); both lie in \(S\),
\(x=v+1\), \(v=v'+1\) and \(\operatorname{Pw}(v)\).  Let \(w,w'\in S\) with
\(\operatorname{Pw}(w)\), \(w=w'+1\) and \(\operatorname{Div}(w',v')\).
Then \(w\geq2\) is a power of \(2\), \(w=2^i\) with \(i\geq1\), and
\(2^i-1\) divides \(2^j-1\); by the divisibility fact \(i\) divides
\(j=2^m\), so \(i=2^a\).  If \(a=0\), \(w=2\) and \(\operatorname{Two}(w)\);
if \(a\geq1\), \(i\) is even and \(\operatorname{Sq}(w)\).  With
\(\operatorname{Prime}(x)\) this gives \(\operatorname{Fer}(x)\).
\end{proof}
\end{annexitem}

\begin{annexitem}{annex:p03-w03-p03-071-prime-graph}{thm:p03-w03-p03-071-prime-graph}{Primality of an encoded length is almost surely first-order}
\begin{proof}
Fix an ordered graph on \(\Omega\) and \(u\in\Omega\), and let
\(c,d,c',d'\) satisfy \(\operatorname{Good}(u,c,d,c',d')\).  By
\cref{thm:p03-w03-p03-071-witness}, item~2, the relations \(P\) and \(T\)
built from these parameters are \(\mathrm{Add}_u\) and \(\mathrm{Mul}_u\) on
\(S=\{1,\ldots,u\}\).  Every quantifier of \(\operatorname{Prime}^{P,T}_u\)
ranges over \(S\); the order of the graph restricted to \(S\) is the order of
\(S\), and for \(x,y\leq u\) the formula \(y=x+1\) of
\cref{def:p03-w03-p03-071-gord} says that \(y\) is the successor of \(x\) in
\(S\), since no element lies strictly between them.  So
\(\operatorname{Prime}^{P,T}_u(u)\) holds in the graph exactly when
\(\operatorname{Prime}(u)\) holds in \((S,<,\mathrm{Add}_u,\mathrm{Mul}_u)\),
that is, by \cref{thm:p03-w03-p03-071-prime}, exactly when \(u\) is prime.
Hence \(\operatorname{Prime}^{\ast}(u)\) holds if and only if
\(\operatorname{Ar}(u)\) holds and \(u\) is prime.  In particular
\(\operatorname{Prime}^{\ast}(u)\) implies that \(u\) is prime, in every
ordered graph.  By \cref{thm:p03-w03-p03-071-representation}, with
probability tending to \(1\), \(\operatorname{Ar}(u)\) holds for every
\(u\leq0.9L\); on that event \(\operatorname{Prime}^{\ast}(u)\) holds exactly
when \(u\) is prime, for all such \(u\).  The argument for
\(\operatorname{Fer}^{\ast}\) is the same, with \(\operatorname{Fer}\) in
place of \(\operatorname{Prime}\).
\end{proof}
\end{annexitem}

\begin{annexitem}{annex:p03-w03-p03-071-chs}{thm:p03-w03-p03-071-chs}{Nonconvergence (Compton, Henson and Shelah)}
\begin{proof}
Let \(E_n\) be the event that \(\operatorname{Ar}(u)\) holds for all \(u\)
with \(1\leq u\leq0.9L\) and fails for all \(u\in\Omega\) with \(u\geq3L\).
By \cref{thm:p03-w03-p03-071-representation,thm:p03-w03-p03-071-upper},
\(\Pr[E_n]\to1\).  Take \(n\) so large that \(L\geq10\),
\(0.9L-1\geq2^{40}\), and \(\lceil3L\rceil\leq n-1\).

Assume \(E_n\).  Let \(W\) be the set of \(u\in\Omega\) with \(u+1\in\Omega\),
\(\operatorname{Ar}(u)\) and \(\neg\operatorname{Ar}(u+1)\).  Since
\(\operatorname{Ar}(1)\), the largest \(u_0\) with \(\operatorname{Ar}(u_0)\)
exists, satisfies \(u_0<3L\leq n-1\), and \(u_0\in W\); so \(W\neq
\varnothing\).  Each \(u\in W\) has \(u<3L\), because
\(\operatorname{Ar}(u)\), and \(u+1>0.9L\), because
\(\neg\operatorname{Ar}(u+1)\); hence \(0.8L\leq0.9L-1<u<3L\) and
\(u\geq2^{40}\).

\(A\) holds exactly when some \(u\in W\) has parameters for which
\[
 \operatorname{Good}(u,c,d,c',d')\quad\text{and}\quad\Phi^{P,T}_u
\]
both hold.  For such parameters \(P=\mathrm{Add}_u\) and \(T=\mathrm{Mul}_u\) on
\(\{1,\ldots,u\}\) by \cref{thm:p03-w03-p03-071-witness}, item~2, and every
quantifier of \(\Phi^{P,T}_u\) ranges over \(\{1,\ldots,u\}\); so
\(\Phi^{P,T}_u\) holds in the graph exactly when \(\Phi\) holds in the
arithmetic structure on \(\{1,\ldots,u\}\) with
\(<,\mathrm{Add}_u,\mathrm{Mul}_u\).  By
\cref{thm:p03-w03-p03-071-log-mod}, this holds exactly when
\(\lfloor\log u\rfloor\bmod40\in\{1,\ldots,20\}\).

Let \(h=h(n)=\lfloor\log(0.8L)\rfloor\).  For \(u\in W\),
\(h\leq\lfloor\log u\rfloor\leq\lfloor\log(3L)\rfloor\leq h+2\), because
\(\log(3L)=\log(0.8L)+\log3.75\) and \(\log3.75<2\).  If \(h\equiv1\pmod
{40}\), every \(u\in W\) has \(\lfloor\log u\rfloor\bmod40\in\{1,2,3\}\),
and \(A\) holds, since \(W\neq\varnothing\).  If \(h\equiv21\pmod{40}\),
every \(u\in W\) has residue in \(\{21,22,23\}\), and \(A\) fails.

The sequence \(h(n)\) is nondecreasing and tends to infinity.  Moreover
\(h(n+1)\leq h(n)+1\) for \(n\geq2\), because \(\log(n+1)\leq8\log n\) gives
\(L(n+1)\leq2L(n)\).  So \(h\) takes every integer value from some point on,
and there are infinitely many \(n\) with \(h(n)\equiv1\) and infinitely many
with \(h(n)\equiv21\pmod{40}\).  Along the first, for large \(n\),
\(\Pr[A]\geq\Pr[E_n]\to1\); along the second,
\(\Pr[A]\leq1-\Pr[E_n]\to0\).  As all values lie in \([0,1]\), the
\(\limsup\) is \(1\) and the \(\liminf\) is \(0\).
\end{proof}
\end{annexitem}

\begin{annexitem}{annex:p03-w03-p03-071-all-sizes}{thm:p03-w03-p03-071-all-sizes}{Common neighbourhoods of every small size}
\begin{proof}
Let \(p=n^{-1/4}\), so \(p^4=\frac1n\), and let \(n\) be large.  Put
\(a=\lceil n^{1/5}\rceil\), \(V_1=\{1,\ldots,a\}\),
\(V_2=\{a+1,\ldots,n\}\), \(b=n-a\), \(K=\lfloor a/4\rfloor\), and
\(Q_j=\{4j+1,4j+2,4j+3,4j+4\}\subseteq V_1\) for \(0\leq j<K\).
For each \(j\) let \(Y_j=|N(Q_j)\cap V_2|\) and
\(Z_j=|N(Q_j)\cap V_1|\), where \(N(Q_j)\) is the common neighbourhood
of the four elements of \(Q_j\).

\emph{Independence of the \(Y_j\).}  A vertex \(w\in V_2\) lies in
\(N(Q_j)\) exactly when the four pairs \(\{q,w\}\), \(q\in Q_j\), are edges.
For fixed \(j\) these events, one for each \(w\in V_2\), use disjoint
sets of pairs and each has probability \(p^4=\frac1n\); so \(Y_j\) has the
binomial distribution with parameters \(b\) and \(\frac1n\).  Distinct
\(j\) use the disjoint pair sets \(Q_j\times V_2\), so
\(Y_0,\ldots,Y_{K-1}\) are independent.

\emph{The probability of size \(s\).}  Let \(0\leq s\leq\sigma(n)\); then
\(s\leq\ln n\).  With \(\pi_s=\Pr[Y_j=s]\),
\[
  \pi_s=\binom bs n^{-s}\Bigl(1-\frac1n\Bigr)^{b-s}
  \geq\frac{(b-s)^s}{s!\,n^s}\Bigl(1-\frac1n\Bigr)^{n}
  \geq\frac1{4\,s!}\Bigl(1-\frac{a+s}n\Bigr)^s ,
\]
because \(b(b-1)\cdots(b-s+1)\geq(b-s)^s\), and \((1-\frac1n)^n\) increases
with \(n\) and equals \(\frac14\) at \(n=2\).  Since \(a+s\leq n^{1/5}+1+\ln
n\leq2n^{1/5}\), Bernoulli's inequality gives
\((1-\frac{a+s}n)^s\geq1-2n^{-4/5}\ln n\geq\frac12\).  So
\(\pi_s\geq1/(8\,s!)\).  For \(s\geq1\), \(\ln s!\leq s\ln s\leq
\sigma(n)\ln\sigma(n)\leq\frac{\ln n}{20\ln\ln n}\cdot\ln\ln n=\frac{\ln
n}{20}\), because \(t\ln t\) increases for \(t\geq1\) and \(\sigma(n)\leq\ln
n\); the same bound holds for \(s=0\).  Hence \(s!\leq n^{1/20}\) and
\(\pi_s\geq n^{-1/20}/8\).

\emph{Many good quadruples.}  Let \(J_s\) be the number of \(j<K\) with
\(Y_j=s\).  By independence \(J_s\) is binomial with parameters \(K\) and
\(\pi_s\), and mean \(\mu_s=K\pi_s\).  As \(K\geq a/4-1\geq n^{1/5}/8\),
\(\mu_s\geq n^{3/20}/64\).  By Chebyshev's inequality,
\[
  \Pr\bigl[J_s\leq\tfrac12\mu_s\bigr]\leq
  \frac{K\pi_s(1-\pi_s)}{(\mu_s/2)^2}\leq\frac4{\mu_s}\leq256\,n^{-3/20}.
\]
Summed over the at most \(\ln n+1\) values of \(s\), the probability that
\(J_s\leq\frac12\mu_s\) for some \(s\leq\sigma(n)\) is at most
\(256(\ln n+1)n^{-3/20}\), which tends to \(0\).

\emph{Few spoiled quadruples.}  For \(w\in V_1\setminus Q_j\),
\(\Pr[w\in N(Q_j)]=\frac1n\), so \(\Pr[Z_j\geq1]\leq\mathbb E[Z_j]\leq
\frac an\).  Let \(D\) be the number of \(j<K\) with \(Z_j\geq1\); then
\(\mathbb E[D]\leq Ka/n\), and by Markov's inequality
\(\Pr[D\geq Kn^{-1/2}]\leq an^{-1/2}\leq2n^{-3/10}\), which tends to \(0\).
\(D\) does not depend on \(s\).

\emph{Conclusion.}  Outside the two exceptional events, for every
\(s\leq\sigma(n)\),
\[
  J_s>\tfrac12\mu_s\geq\frac{K\,n^{-1/20}}{16}>K\,n^{-1/2}>D,
\]
the middle strict inequality holding once \(n^{9/20}>16\).  So some
\(j<K\) has \(Y_j=s\) and \(Z_j=0\).  Every common neighbour of \(Q_j\) lies
in \(V_1\) or in \(V_2\), hence \(|N(Q_j)|=Y_j+Z_j=s\), and the
four elements of \(Q_j\) are the required vertices.
\end{proof}
\end{annexitem}

\begin{annexitem}{annex:p03-w03-p03-071-coupling}{thm:p03-w03-p03-071-coupling}{The coupling identity}
\begin{proof}
Fix a \(k\)-subset \(S\) with elements \(t_1<\cdots<t_k\).  The order
isomorphism \(i\mapsto t_i\) is an isomorphism from the transported graph
on \(\{1,\ldots,k\}\) onto \(G|S\), in which \(i\sim j\) exactly when
\(t_i\sim t_j\) in \(G\).  These \(\binom k2\) adjacencies are distinct pairs
of \(G\), so they are independent, each with probability \(p\).  So, given
\(S\), the transported graph is distributed as \(G_<(k,p)\); as \(S\) is
independent of \(G\), the same holds unconditionally.  Satisfaction of a
sentence is invariant under isomorphism, so
\(f_A(k)=\sum_G\mu(G)\Pr[G^{(k)}\entails A]\) for \(k=n\) and \(k=n+1\), with
\(N=2n+1\).  Subtracting gives the identity.
\end{proof}
\end{annexitem}

\begin{annexitem}{annex:p03-w03-p03-071-restriction-facts}{thm:p03-w03-p03-071-restriction-facts}{Three facts about restrictions}
\begin{proof}
For a coordinate \(i\) and an input \(x\), write \(x^{\oplus i}\) for \(x\)
with coordinate \(i\) flipped.  The probability that a \(\delta\)-random
restriction of \(V\) equals a given \((J\mid z)\) is
\[
  P_\delta(J\mid z)=\delta^{|J|}\Bigl(\frac{1-\delta}2\Bigr)^{|V\setminus J|},
\]
because \(J\) has probability \(\delta^{|J|}(1-\delta)^{|V\setminus J|}\) and
then \(z\) has probability \(2^{-|V\setminus J|}\).

\emph{Item 1.}  \(\mathbf I[f_{J\mid z}]=\sum_{i\in V}\mathbf 1_{i\in
J}\operatorname{Inf}_i[f_{J\mid z}]\).  Fix \(i\in V\) and a set \(J\ni i\).
For \(y\) uniform on \(\{0,1\}^J\) and \(z\) uniform on
\(\{0,1\}^{V\setminus J}\), independent, the combined input \(x=(y,z)\) is
uniform on \(\{0,1\}^V\), and \(f_{J\mid z}(y)\neq f_{J\mid z}(y^{\oplus i})\)
exactly when \(f(x)\neq f(x^{\oplus i})\).  Averaging over \(z\),
\(\mathbb E_z[\operatorname{Inf}_i[f_{J\mid z}]]=\Pr_x[f(x)\neq
f(x^{\oplus i})]=\operatorname{Inf}_i[f]\), whatever \(J\ni i\) is.  As
\(\Pr[i\in J]=\delta\), \(\mathbb E[\mathbf 1_{i\in J}\operatorname{Inf}_i[f_{J\mid
z}]]=\delta\operatorname{Inf}_i[f]\); summing over \(i\) gives item~1.

\emph{Item 2.}  Fix a restriction \((J'\mid y)\) of \(V\).  The composed
restriction equals it exactly when the first restriction is
\((J\mid y|_{V\setminus J})\) for some \(J\) with \(J'\subseteq J\) and the
second is \((J'\mid y|_{J\setminus J'})\).  Summing over \(J\), with
\(D=V\setminus J'\) and \(H=J\setminus J'\subseteq D\),
\begin{align*}
  \Pr[\,\cdot\,]&=\sum_{H\subseteq D}\delta_1^{|J'|+|H|}
     \Bigl(\frac{1-\delta_1}2\Bigr)^{|D|-|H|}\,
     \delta_2^{|J'|}\Bigl(\frac{1-\delta_2}2\Bigr)^{|H|}\\
  &=(\delta_1\delta_2)^{|J'|}2^{-|D|}\sum_{H\subseteq D}
     \bigl(\delta_1(1-\delta_2)\bigr)^{|H|}(1-\delta_1)^{|D|-|H|}\\
   &=(\delta_1\delta_2)^{|J'|}\Bigl(\frac{1-\delta_1\delta_2}2\Bigr)^{|D|},
\end{align*}
by the binomial theorem, since \(\delta_1(1-\delta_2)+1-\delta_1=
1-\delta_1\delta_2\).  This is \(P_{\delta_1\delta_2}(J'\mid y)\).

\emph{Item 3.}  Let a decision tree of depth at most \(k\) compute \(g\).  On
input \(y\) the tree queries at most \(k\) coordinates.  If coordinate \(i\)
is not queried on the path of \(y\), then \(y^{\oplus i}\) follows the same
path, which depends only on the queried values, and reaches the same leaf;
so \(i\) is not pivotal at \(y\).  Hence at most \(k\) coordinates are
pivotal at each \(y\), and \(\mathbf I[g]=\sum_i\Pr_y[i\text{ pivotal at
}y]=\mathbb E_y[\#\{i:i\text{ pivotal at }y\}]\leq k\).  For each leaf
\(\lambda\), let \(\tau_\lambda\) be the AND of the literals \(s_i\) or
\(\neg s_i\) recording the edges on the path to \(\lambda\); its coordinates
are distinct, and \(y\) satisfies \(\tau_\lambda\) exactly when \(y\)
reaches \(\lambda\).  Each \(y\) reaches exactly one leaf, so \(g\) is the OR
of \(\tau_\lambda\) over the leaves labelled \(1\), a DNF of width at most
\(k\), and the AND of \(\neg\tau_\lambda\) over the leaves labelled \(0\),
which by De Morgan's law is a CNF of width at most \(k\).  Finally each
\(\operatorname{Inf}_i[g]\) is a probability, so \(\mathbf I[g]\leq|J|\).
\end{proof}
\end{annexitem}

\begin{annexitem}{annex:p03-w03-p03-071-switching}{thm:p03-w03-p03-071-switching}{Switching lemma (H\aa stad)}
\begin{proof}
\emph{Reduction to DNFs.}  If \(F\) is a CNF of width at most \(w\), its
negation is, by De Morgan's law, a DNF of width at most \(w\), and
\((\neg F)_{J\mid z}=\neg(F_{J\mid z})\); a decision tree computes a function
exactly when the same tree with leaf labels exchanged computes its negation,
so the two have the same \(\operatorname{DT}\).  So let \(F=T_1\vee\cdots\vee
T_M\) be a DNF of width at most \(w\), with its terms in this fixed order,
the literals of each term listed in a fixed order, and \(V\) listed in a
fixed order.  If \(w=0\), every term is empty, \(F\) is constant and the
claim is trivial; so let \(w\geq1\).  Write \(q=\frac{1-\delta}2\geq\frac14\).

\emph{Terminology.}  A restriction \(\rho'\) \emph{extends} \(\rho\) if every
coordinate fixed by \(\rho\) is fixed by \(\rho'\) to the same value.  For
\(\rho=(J\mid z)\) and an assignment \(\pi\in\{0,1\}^S\) with \(S\subseteq
J\), \(\rho\pi\) is the restriction \((J\setminus S\mid(z,\pi))\).  A term is
\emph{falsified} by \(\rho\) if one of its literals is fixed false, and
\emph{satisfied} if all its literals are fixed true.  Both properties pass to
extensions.

\emph{The canonical decision tree.}  For a restriction \(\rho=(J\mid z)\)
define a tree \(\operatorname{CDT}(\rho)\) by recursion on \(|J|\).  If
\(F_\rho\) is constant, \(\operatorname{CDT}(\rho)\) is a leaf with that
value.  Otherwise some term is not falsified by \(\rho\), since otherwise
\(F_\rho\equiv0\); let \(T\) be the first such term.  \(T\) is not satisfied,
since otherwise \(F_\rho\equiv1\); so the set \(S\subseteq J\) of
coordinates of \(T\) that are free under \(\rho\) is nonempty, and the other
literals of \(T\) are fixed true.  Then \(\operatorname{CDT}(\rho)\) queries
the coordinates of \(S\) one after another in the order of \(V\), forming a
complete binary tree of depth \(|S|\), and below the node reached by
\(\pi\in\{0,1\}^S\) it continues with \(\operatorname{CDT}(\rho\pi)\), where
\(|J\setminus S|<|J|\).  By induction on \(|J|\), \(\operatorname{CDT}(\rho)\)
computes \(F_\rho\): an input \(y\in\{0,1\}^J\) reaches the subtree for
\(\pi=y|_S\), which computes \(F_{\rho\pi}(y|_{J\setminus S})=F_\rho(y)\).
No coordinate is queried twice on a path, since queried coordinates are
fixed in the later restrictions.  So \(\operatorname{DT}(F_\rho)\) is at
most the depth of \(\operatorname{CDT}(\rho)\).

A root-to-leaf path of \(\operatorname{CDT}(\rho)\) therefore consists of
blocks
\[
 (T^{(1)},S_1,\pi_1),\ldots,(T^{(l)},S_l,\pi_l).
\]
Here \(T^{(j)}\) is the first term not falsified by \(\rho\pi_1\cdots\pi_{j-1}\), \(S_j\) is its set
of free coordinates under that restriction, and \(\pi_j\in\{0,1\}^{S_j}\) is
the part of the path inside the block.  The sets \(S_j\) are pairwise
disjoint, the coordinates of \(T^{(j)}\) outside \(S_j\) are fixed true by
\(\rho\pi_1\cdots\pi_{j-1}\), and the length of the path is
\(\sum_j|S_j|\).

\emph{The encoding.}  Fix \(k\geq1\) and let \(B\) be the set of
restrictions \(\rho\) for which \(\operatorname{CDT}(\rho)\) has depth at
least \(k\).  For \(\rho\in B\) choose a root-to-leaf path of length at least
\(k\), with blocks as above, let \(r\) be the least index with
\(|S_1|+\cdots+|S_r|\geq k\), and let \(S'_r\) consist of the first
\(k-(|S_1|+\cdots+|S_{r-1}|)\) coordinates of \(S_r\) in the query order.
Put \(S'_j=S_j\) for \(j<r\).  Thus \(S'_1,\ldots,S'_r\) are disjoint, free
under \(\rho\), and have \(k\) elements together.  For \(j\leq r\) let
\(\sigma_j\in\{0,1\}^{S'_j}\) give each coordinate the value that makes its
literal in \(T^{(j)}\) true, and put
\[
  \rho^\ast=\rho\sigma_1\cdots\sigma_r .
\]
The \emph{record} of \(\rho\) lists, for each of the \(k\) coordinates of
\(S'_1,\ldots,S'_r\) in order, three items: the position (in
\(\{1,\ldots,w\}\)) of its literal in the term \(T^{(j)}\) of its block; a
bit telling whether it is the last coordinate of \(S'_j\); and its value
under \(\pi_j\).  There are at most \((4w)^k\) records.

\emph{Decoding.}  We show that \(\rho\) is determined by \(\rho^\ast\) and
its record.  For \(1\leq j\leq r\) let
\(\rho_j=\rho\pi'_1\cdots\pi'_{j-1}\sigma_j\cdots\sigma_r\), where
\(\pi'_i=\pi_i|_{S'_i}\); so \(\rho_1=\rho^\ast\).  \emph{Claim:}
\(T^{(j)}\) is the first term not falsified by \(\rho_j\).  Since
\(\pi'_i=\pi_i\) for \(i<r\), \(\rho_j\) extends
\(\rho\pi_1\cdots\pi_{j-1}\), which falsifies every term before
\(T^{(j)}\).  The coordinates of \(T^{(j)}\) outside \(S_j\)
are fixed true by \(\rho\pi_1\cdots\pi_{j-1}\), and \(\rho_j\) agrees with it
there; on \(S'_j\), \(\rho_j\) is \(\sigma_j\), which makes the literals
true; the blocks \(S'_i\) with \(i>j\) are disjoint from the coordinates of
\(T^{(j)}\), all of which are fixed by \(\rho\pi_1\cdots\pi_j\) while
\(S'_i\) is free there; and the coordinates of \(S_r\setminus S'_r\), when
\(j=r\), are free under \(\rho_r\).  So no literal of \(T^{(j)}\) is fixed
false by \(\rho_j\), which proves the claim.

The decoder proceeds for \(j=1,\ldots,r\).  Knowing \(\rho_j\), it finds
\(T^{(j)}\) as the first term not falsified by \(\rho_j\).  The record,
read from where block \(j-1\) ended up to the next ``last'' bit, gives the
positions in \(T^{(j)}\) of the coordinates of \(S'_j\), hence \(S'_j\)
itself, and their values under \(\pi_j\); replacing \(\sigma_j\) by
\(\pi'_j\) on \(S'_j\) turns \(\rho_j\) into \(\rho_{j+1}\).  The number
\(r\) of blocks is the number of ``last'' bits in the record.  At the end
the decoder knows \(S'_1\cup\cdots\cup S'_r\); making these coordinates
free in \(\rho^\ast\) returns \(\rho\).  Hence the map
\(\rho\mapsto(\rho^\ast,\text{record})\) is injective on \(B\).

\emph{Counting.}  \(\rho^\ast\) fixes \(k\) more coordinates than \(\rho\),
so \(P_\delta(\rho)=(\delta/q)^kP_\delta(\rho^\ast)\), with \(P_\delta\) as in
the proof of \cref{thm:p03-w03-p03-071-restriction-facts}.  By injectivity
each restriction occurs as \(\rho^\ast\) for at most \((4w)^k\) members of
\(B\).  So
\begin{align*}
  \Pr[B]&=\sum_{\rho\in B}P_\delta(\rho)
  =\Bigl(\frac\delta q\Bigr)^k\sum_{\rho\in B}P_\delta(\rho^\ast)\\
  &\leq\Bigl(\frac\delta q\Bigr)^k(4w)^k\sum_{\rho'}P_\delta(\rho')
  =\Bigl(\frac{4w\delta}q\Bigr)^k\leq(16\delta w)^k,
\end{align*}
the sum over all restrictions \(\rho'\) of \(V\) being \(1\), and \(q\geq
\frac14\).  Since \(\operatorname{DT}(F_{J\mid z})\geq k\) implies
\((J\mid z)\in B\), the lemma follows.
\end{proof}
\end{annexitem}

\begin{annexitem}{annex:p03-w03-p03-071-ac0-influence}{thm:p03-w03-p03-071-ac0-influence}{Total influence of bounded-depth circuits}
\begin{proof}
For each node \(v\) of the circuit let \(g_v:\{0,1\}^N\to\{0,1\}\) be the
function it computes.  By the convention of
\cref{def:p03-w03-p03-071-circuit}, a layer-\(1\) node is joined to each
variable at most once, so it computes a term or a clause; we may assume the
layer-\(1\) nodes are AND gates, the other case being dual (exchange AND
with OR, terms with clauses and DNFs with CNFs throughout).  Since
\(\log_2\frac43>0.415\) and \(3\cdot0.415>1\),
\((\frac43)^t\geq(2Ns)^{1.245}\geq Ns\), so \(Ns(\frac34)^t\leq1\).

\emph{The restrictions.}  Let \(\rho_0\) be a \(\frac12\)-random restriction
of \(\{1,\ldots,N\}\) and, for \(1\leq j\leq d-1\), given
\(\rho_0,\ldots,\rho_{j-1}\), let \(\rho_j\) be a \(\frac1{32t}\)-random
restriction of the coordinates still free.  Let \(R_j\) be the composition
of \(\rho_0,\ldots,\rho_{j-1}\), and \(\rho=R_d\).  By
\cref{thm:p03-w03-p03-071-restriction-facts}, item~2, applied \(d-1\) times,
\(\rho\) is a \(\delta\)-random restriction with
\(\delta=\frac12(32t)^{-(d-1)}\).  Call a node \(v\) in layer \(j\),
\(1\leq j\leq d-1\), \emph{switched} if
\(\operatorname{DT}((g_v)_{R_j})\leq t-1\), and let \(\mathcal G\) be the
event that every node in layers \(1\) to \(d-1\) is switched.

\emph{Layer \(1\).}  Let \(v\) be a layer-\(1\) AND gate with \(m\)
literals.  If \(m\leq t-1\), then \((g_v)_{R_1}\) is a constant or an AND of
at most \(t-1\) free literals, which a decision tree of that depth computes;
so \(v\) is switched.  If \(m\geq t\): each literal is fixed false by
\(\rho_0\) with probability \(\frac14\), independently since the
coordinates are distinct, and if one is, \((g_v)_{R_1}\equiv0\).  So \(v\)
fails to be switched with probability at most \((\frac34)^m\leq(\frac34)^t\).

\emph{Layer \(j+1\), for \(1\leq j\leq d-2\).}  Fix \(R_j\) such that every
node in layer \(j\) is switched, and let \(v\) be a node in layer \(j+1\),
say an OR of nodes \(v_1,\ldots,v_a\) in layer \(j\) (an AND is dual).  By
\cref{thm:p03-w03-p03-071-restriction-facts}, item~3, each
\((g_{v_i})_{R_j}\) is a DNF of width at most \(t-1\) over the free
coordinates, so \((g_v)_{R_j}\), their OR, is one too.  As
\((g_v)_{R_{j+1}}=((g_v)_{R_j})_{\rho_j}\),
\cref{thm:p03-w03-p03-071-switching} with \(\delta=\frac1{32t}\leq\frac12\),
\(w=t-1\) and \(k=t\) shows that, conditionally on \(R_j\), \(v\) fails to
be switched with probability at most \((16(t-1)/(32t))^t\leq2^{-t}\leq
(\frac34)^t\).

Summing over the \(s\) nodes of layers \(1\) to \(d-1\): \(\mathcal G\) can
fail only if, for some \(j\leq d-1\), every node of layer \(j-1\) is switched
(no condition for \(j=1\)) and some node of layer \(j\) is not.  So
\(\Pr[\neg\mathcal G]\leq s(\frac34)^t\).

\emph{The output.}  Fix \(R_{d-1}\) in \(\mathcal G\).  As in the previous
step, \((g_{\mathrm{out}})_{R_{d-1}}=f_{R_{d-1}}\) is a DNF or CNF of width at
most \(t-1\), and by \cref{thm:p03-w03-p03-071-switching},
\[
  \mathbb E\bigl[\operatorname{DT}(f_\rho)\bigm|R_{d-1}\bigr]
  =\sum_{r\geq1}\Pr\bigl[\operatorname{DT}(f_\rho)\geq r\bigm|R_{d-1}\bigr]
  \leq\sum_{r\geq1}2^{-r}=1 .
\]

\emph{Conclusion.}  By \cref{thm:p03-w03-p03-071-restriction-facts},
item~3, \(\mathbf I[f_\rho]\leq\operatorname{DT}(f_\rho)\), and always
\(\mathbf I[f_\rho]\leq N\).  The event \(\mathcal G\) depends only on
\(R_{d-1}\), so
\[
  \mathbb E\bigl[\mathbf I[f_\rho]\bigr]\leq
  \mathbb E\bigl[\mathbf 1_{\mathcal G}\,\mathbb E[\operatorname{DT}(f_\rho)
  \mid R_{d-1}]\bigr]+N\Pr[\neg\mathcal G]\leq1+Ns\bigl(\tfrac34\bigr)^t
  \leq2 .
\]
By item~1 of \cref{thm:p03-w03-p03-071-restriction-facts},
\(\delta\,\mathbf I[f]=\mathbb E[\mathbf I[f_\rho]]\leq2\), so
\(\mathbf I[f]\leq2/\delta=4(32t)^{d-1}\).  Finally
\(t\leq3\log(2Ns)+1\leq4\log(2Ns)\), as \(\log(2Ns)\geq1\), so
\(32t\leq128\log(2Ns)\).
\end{proof}
\end{annexitem}

\begin{annexitem}{annex:p03-w03-p03-071-translation}{thm:p03-w03-p03-071-translation}{A first-order sentence becomes a bounded-depth circuit}
\begin{proof}
Abbreviations \(\to,\leftrightarrow\) are unfolded first, and \(\ell\) counts
the result.  Fix \(G\).  A \emph{formula tree} is a finite rooted tree whose
leaves are labelled by literals \(s_b,\neg s_b\) or constants \(0,1\) and
whose internal nodes are labelled AND or OR; it computes a function of
\(s\in\{0,1\}^N\) in the evident way.

\emph{Construction.}  For a formula \(\psi(x_1,\ldots,x_k)\) and
\(\bar a\in\{1,\ldots,N\}^k\) define two trees \(C^+(\psi,\bar a)\) and
\(C^-(\psi,\bar a)\) by recursion on \(\psi\).  For atomic \(\psi\),
\(C^+\) is the constant leaf giving the truth value of \(\psi[\bar a]\) in
\(G\) (each of \(a_i\sim a_j\), \(a_i<a_j\), \(a_i=a_j\) is fixed once \(G\)
and \(\bar a\) are), and \(C^-\) the opposite constant.  Further
\(C^\pm(\neg\psi)=C^\mp(\psi)\); \(C^+(\psi\wedge\chi)\) is an AND of
\(C^+(\psi),C^+(\chi)\) and \(C^-(\psi\wedge\chi)\) an OR of
\(C^-(\psi),C^-(\chi)\); dually for \(\vee\); and
\begin{align*}
  C^+(\exists_x\psi,\bar a)&=\mathrm{OR}_{b=1}^{N}\,\mathrm{AND}\bigl(s_b,
     C^+(\psi,\bar ab)\bigr),\\
  C^-(\exists_x\psi,\bar a)&=\mathrm{AND}_{b=1}^{N}\,\mathrm{OR}\bigl(\neg s_b,
     C^-(\psi,\bar ab)\bigr),
\end{align*}
and dually for \(\forall_x\).

\emph{Correctness.}  By induction on \(\psi\): for every \(S\) with indicator
\(s\) and every \(\bar a\in S^k\), \(C^+(\psi,\bar a)(s)=1\) exactly when
\(G|S\entails\psi[\bar a]\), and \(C^-(\psi,\bar a)(s)=1\) exactly when not.
Atomic formulas are evaluated in \(G|S\) as in \(G\), because \(G|S\) is
induced.  The connective steps are immediate.  For \(\exists_x\psi\):
\(C^+(s)=1\) iff some \(b\) has \(s_b=1\) and \(C^+(\psi,\bar ab)(s)=1\),
iff (induction, as \(\bar ab\in S^{k+1}\)) some \(b\in S\) has
\(G|S\entails\psi[\bar ab]\); and \(C^-(s)=1\) iff every \(b\in S\) has
\(G|S\not\entails\psi[\bar ab]\).  The \(\forall\) step is dual.  For the
sentence \(A\), \(f_G\) is computed by \(C^+(A)\).

\emph{Size and height.}  Let \(\lvert\psi\rvert\) be the number of
occurrences in \(\psi\) and \(\operatorname{qr}(\psi)\) its quantifier rank.
By induction, \(C^\pm(\psi,\bar a)\) has at most
\(\lvert\psi\rvert(3N)^{\operatorname{qr}(\psi)}\) nodes and height at most
\(2\lvert\psi\rvert\).  Atomic: one node, height \(0\).  Negation changes
neither.  Binary connectives add one node and one level, and
\(\lvert\psi\wedge\chi\rvert=\lvert\psi\rvert+\lvert\chi\rvert+1\).  For
\(\exists_x\psi\) of rank \(q\), the tree has
\(1+N(2+X)\) nodes with \(X\leq\lvert\psi\rvert(3N)^{q-1}\), and
\(1+2N+N\lvert\psi\rvert(3N)^{q-1}\leq(3N)^q+(3N)^q\lvert\psi\rvert=
(3N)^q\lvert\exists_x\psi\rvert\); the height grows by \(2\).  Hence
\(C^+(A)\) has at most \(\ell(3N)^q\) nodes and height at most \(2\ell\).

\emph{Normalization.}  Apply the following steps while one applies; each
removes at least one node and none increases the height or changes the
function computed: (a) an AND with a child \(0\), or an OR with a child
\(1\), becomes that constant; (b) a child \(1\) of an AND, or \(0\) of an OR,
is deleted; (c) an AND without children becomes \(1\), an OR without
children \(0\); (d) an AND with children \(s_b\) and \(\neg s_b\) becomes
\(0\), an OR with both becomes \(1\); (e) a repeated literal child is
deleted; (f) a gate child with the label of its parent is removed and its
children attached to the parent.  At the end the tree is a constant, or a
single literal, or every gate has at least one child, no child is a
constant, no gate has a repeated or complementary pair of literal children,
and the labels alternate along every root-to-leaf path.

\emph{Layering.}  In the last case let \(h\) be the height and
\(d=\max(2,h)\leq d_A\).  If \(h=1\), add a new root with one child, the old
root, and the other label; then \(h=2\).  A gate at distance \(t\) from the
root is put in layer \(d-t\); a gate has a child, so \(t\leq d-1\), and
gates in the same layer carry the same label, the labels alternating from
layer to layer.  A literal leaf at distance \(t<d\) from the root, whose
parent therefore lies in layer \(d-t+1\), is replaced by a chain of \(d-t\)
single-child gates placed in layers \(d-t,\ldots,1\), each carrying the
label of its layer, the last of which has the literal as its only child; a
leaf at distance \(d\) is left as it is.  Single-child gates compute their
child.  Every child of the root is now a gate in layer \(d-1\geq1\), so the
size is at least \(1\).  Identifying all leaves with the same literal with one layer-\(0\)
node, and adding the unused literals there, gives a depth-\(d\) circuit in
the sense of \cref{def:p03-w03-p03-071-circuit} computing \(f_G\).  A
layer-\(1\) node is either an original gate, whose children are distinct and
non-complementary literals, or a chain end with one literal child.  The
nodes in layers \(1\) to \(d-1\) are original gates, at most \(\ell(3N)^q\),
the added root, and chain gates, at most \(d-1\leq d_A\) per leaf, with at
most \(\ell(3N)^q\) leaves; so the size is at most
\((d_A+1)\ell(3N)^q+2\leq s_N\).
\end{proof}
\end{annexitem}

\begin{annexitem}{annex:p03-w03-p03-071-slice}{thm:p03-w03-p03-071-slice}{Adjacent slices}
\begin{proof}
Let \(E_k\) be the set of pairs \((x,b)\) with \(x\) having exactly \(k\)
ones and \(x_b=0\); \(|E_k|=\binom Nk(N-k)\).  Choose \((x,b)\) uniformly
from \(E_k\) and let \(y=x+e_b\), where \(e_b\) is the \(b\)-th unit vector.
Each \(x\) of weight \(k\) occurs in exactly \(N-k\) pairs, so \(x\) is
uniform on weight \(k\).  Each \(y\) of weight \(k+1\) arises from exactly
\(k+1\) pairs, one for each \(b\) with \(y_b=1\), so \(y\) is uniform on
weight \(k+1\).  Hence
\(\mu_{k+1}(f)-\mu_k(f)=\mathbb E[f(y)-f(x)]\), and
\(|\mu_{k+1}(f)-\mu_k(f)|\leq\Pr[f(x)\neq f(y)]\).

Call an edge \(\{x,x+e_b\}\) of the cube a boundary edge if \(f\) differs at
its ends.  For each \(b\), coordinate \(b\) is pivotal at \(s\) exactly when
the edge in direction \(b\) at \(s\) is a boundary edge, and each edge has
two ends, so \(\operatorname{Inf}_b[f]=2B_b/2^N\) with \(B_b\) the number of
boundary edges in direction \(b\).  Summing, \(\mathbf I[f]=B/2^{N-1}\),
where \(B\) is the number of all boundary edges.  Distinct pairs of \(E_k\)
are distinct edges, so \(\Pr[f(x)\neq f(y)]\leq B/|E_k|=2^{N-1}\mathbf I[f]/
(\binom Nk(N-k))\).

For \(N=2n+1\), \(k=n\): \(\binom{2n+1}{n}(n+1)=(2n+1)\binom{2n}{n}\) and
\(2^{N-1}=4^n\).  Let \(\pi_n=\binom{2n}n4^{-n}=\prod_{j=1}^n\frac{2j-1}{2j}\).
Since \((2j-1)^2\geq(2j-2)(2j)\), we have
\(\bigl(\frac{2j-1}{2j}\bigr)^2\geq\frac{j-1}{j}\) for \(j\geq2\), so
\(\pi_n^2\geq\frac14\prod_{j=2}^n\frac{j-1}{j}=\frac1{4n}\) and
\(\pi_n\geq\frac1{2\sqrt n}\).  Hence the bound is at most
\(\mathbf I[f]\cdot2\sqrt n/(2n+1)<\mathbf I[f]/\sqrt n\).
\end{proof}
\end{annexitem}

\begin{annexitem}{annex:p03-w03-p03-071-advantage}{thm:p03-w03-p03-071-advantage}{Bounded-depth circuits cannot tell \(n\) true inputs from \(n+1\)}
\begin{proof}
Let \(N=2n+1\).  By \cref{thm:p03-w03-p03-071-slice},
\(|\mu_{n+1}(f)-\mu_n(f)|\leq\mathbf I[f]/\sqrt n\).  By
\cref{thm:p03-w03-p03-071-ac0-influence},
\(\mathbf I[f]\leq4(128\log(2Ns))^{d-1}\), and \(2Ns\leq2N^{c+1}\leq
(2N)^{c+1}\), so \(\log(2Ns)\leq(c+1)\log(2N)=(c+1)\log(4n+2)\).  This gives
the bound.  For fixed \(d\) and \(c\) its numerator is a power of
\(\log(4n+2)\), and \((\log(4n+2))^{d-1}/\sqrt n\to0\) because
\(\log(4n+2)\leq\log(6n)\leq6\log n\) for \(n\geq2\), and for a fixed
integer \(a\geq0\), writing \(\sqrt n=\exp(y)\) with \(y=\frac12\ln n\), the
exponential series gives \(\exp(y)\geq y^{a+1}/(a+1)!\), so
\((\log n)^a/\sqrt n=(2y/\ln2)^a/\exp(y)\leq(2/\ln2)^a(a+1)!/y\to0\).
\end{proof}
\end{annexitem}

\begin{annexitem}{annex:p03-w03-p03-071-shelah}{thm:p03-w03-p03-071-shelah}{Slow oscillation (Shelah)}
\begin{proof}
Fix \(A\), \(n\geq2\), \(N=2n+1\) and \(G\) on \(\{1,\ldots,N\}\).  A
uniform \(k\)-subset \(S\) corresponds to a uniform input of weight \(k\),
so \(\Pr[G^{(k)}\entails A]=\mu_k(f_G)\) with \(f_G\) as in
\cref{thm:p03-w03-p03-071-translation}.  By
\cref{thm:p03-w03-p03-071-slice},
\[
  \bigl|\Pr[G^{(n+1)}\entails A]-\Pr[G^{(n)}\entails A]\bigr|\leq
  \mathbf I[f_G]/\sqrt n .
\]
By \cref{thm:p03-w03-p03-071-translation} there are three cases.  A
constant \(f_G\) has \(\mathbf I=0\), and a literal has \(\mathbf I=1\).
Otherwise \(f_G\) is computed by a depth-\(d\) circuit with
\(2\leq d\leq d_A\) and size \(s\) with \(1\leq s\leq s_N\); the size is at
least \(1\) by the construction in that proof.  Then
\cref{thm:p03-w03-p03-071-ac0-influence} gives
\[
  \mathbf I[f_G]\leq4\bigl(128\log(2Ns)\bigr)^{d-1}
  \leq4\bigl(128\log(2Ns_N)\bigr)^{d_A-1},
\]
because the base is at least \(128\), \(s\leq s_N\) and \(d\leq d_A\).
Since \(N=2n+1\leq3n\) and \(s_N=(d_A+1)\ell(3N)^q+2\leq
(d_A+1)(\ell+2)(3N)^q\), we have \(2Ns_N\leq(d_A+1)(\ell+2)(9n)^{q+1}\); as
\(\log9<4\) and \(\log n\geq1\) for \(n\geq2\),
\(\log(2Ns_N)\leq C_A\log n\) with
\(C_A=\log\bigl((d_A+1)(\ell+2)\bigr)+5(q+1)\geq1\), which depends only on
\(A\).  Hence in every case
\(\mathbf I[f_G]\leq K(\log n)^{d_A-1}\) with \(K=4(128C_A)^{d_A-1}\), using
\(\log n\geq1\).  This gives the displayed bound, uniformly in \(G\).

For the limit, by \cref{thm:p03-w03-p03-071-coupling},
\[
 |f_A(n+1)-f_A(n)|\leq\sum_G\mu(G)\,\frac{K(\log n)^{d_A-1}}{\sqrt n}
 =\frac{K(\log n)^{d_A-1}}{\sqrt n},
\]
which tends to \(0\).
\end{proof}
\end{annexitem}

\begin{annexitem}{annex:p03-w03-p03-071-evenness}{thm:p03-w03-p03-071-evenness}{No almost sure representation of evenness}
\begin{proof}
If \(f_A(2n)\to1\) and \(f_A(2n+1)\to0\), then
\(f_A(2n+1)-f_A(2n)\to-1\).  This is a subsequence of
\(f_A(m+1)-f_A(m)\), which tends to \(0\) by
\cref{thm:p03-w03-p03-071-shelah}; contradiction.
\end{proof}
\end{annexitem}

\begin{annexitem}{annex:p03-w03-p03-071-mso}{thm:p03-w03-p03-071-mso}{The first-order restriction is necessary}
\begin{proof}
Let the vertices be \(1<\cdots<n\).  Suppose \(X\) satisfies the matrix.
Then \(1\in X\), and by induction on \(j\), since \(j+1\) is the successor
of \(j\), \(j\in X\) exactly when \(j\) is odd.  The last clause says
\(n\notin X\), so \(n\) is even.  Conversely, if \(n\) is even, the set of
odd vertices satisfies all three clauses.  The sentence does not mention
\(\sim\), so \(\Pr[G_<(m,p)\entails\operatorname{EV}]\) is \(1\) for even
\(m\) and \(0\) for odd \(m\), for every \(p\).  Then
\(f_{\operatorname{EV}}(2n+1)-f_{\operatorname{EV}}(2n)=-1\) for all \(n\),
so the conclusions of \cref{thm:p03-w03-p03-071-shelah,thm:p03-w03-p03-071-evenness}
fail for \(\operatorname{EV}\).
\end{proof}
\end{annexitem}
